\section{Related Work}
\label{sec:related}

\subsection{Dynamic Inference in Object Detection}
\label{sec:related_dynamic}

Adaptive inference dynamically allocates computation based on the input rather than applying a fixed static cost to every sample. 
While extensively explored in image classification---via early-exit networks~\cite{teera2016branchynet,huang2018multiscale,yang2020ranet,ilhan2024eenet}, spatially adaptive computation~\cite{figurnov2017spatially}, and slimmable or skippable architectures~\cite{yu2018slimmable,yu2019universally,cai2020ofa,wang2018skipnet,wu2018blockdrop,kang2024adaptive}---this paradigm has only recently been extended to the more complex domain of object detection. 
Recent efforts include early-exit detectors~\cite{yang2024adadet,kuhse2025anytimeyolo}, multi-path subnetworks~\cite{heo2020realtime}, deadline-aware LiDAR detectors~\cite{soyyigit2024valo}, and cascaded systems like DynamicDet~\cite{lin2023dynamicdet}. 

We build upon AnyDepth-YOLO~\cite{kang2026anydepth}, which exposes multiple execution depths from a single set of weights by selectively bypassing refinement stages. 
Unlike prior works that focus heavily on the architectural design of these conditional models, we focus on the orthogonal and critical challenge: \emph{how to optimally govern the routing decision at inference time}.



\subsection{Learned Routing for Dynamic Inference}
\label{sec:related_routing}

Many dynamic-inference methods—including SkipNet~\cite{wang2018skipnet}, BlockDrop~\cite{wu2018blockdrop}, AR-Net~\cite{meng2020ar}, and cost-aware NAS~\cite{dai2020danas}—formulate routing as a discrete policy optimized via a weighted sum of accuracy and computation. This policy-centric template carries two fundamental structural flaws. 
First, a fixed computation weight ties the learned router to a single operating point, meaning adaptation to a new budget inevitably requires retraining. 
Second, discrete routing policies inherently suffer from mode collapse, heavily favoring a single path and necessitating fragile workarounds like auxiliary load-balancing losses~\cite{shazeer2017,fedus2022switch,riquelme2021scaling}.
Even DynamicDet~\cite{lin2023dynamicdet}, the closest work to ours, fails to fully escape this architectural trap. 
While it circumvents the first flaw by employing a settable threshold to expose a continuous trade-off without retraining, it still optimizes a routing policy against difficulty-weighted detection losses.
Consequently, the objective predictably collapses onto the more accurate branch, forcing the introduction of a dedicated adaptive-offset mechanism to artificially balance the routes. 
Furthermore, by cascading two independent models, it inherently doubles the parameter storage on the deeper route.

In contrast, our work formulates routing as a supervised advantage-regression
problem. Rather than optimizing a routing policy under an efficiency-aware
objective, the router directly predicts the expected benefit of deeper
execution from frozen detector features. This formulation decouples router
training from deployment-time budget selection, allowing a single trained
router to support a wide range of operating points through threshold
adjustment alone.

\subsection{Threshold-based Decision}
